
\subsubsection*{3.3.3}
\phantomsection
\label{sga3:VIIB:3.3.3}

Continue to suppose that \(k\) is a field of characteristic \(0\).
Let \(\overline k\) be an algebraic closure of \(k\), and let
\(\Gamma\) be the topological Galois group of \(\overline k\) over
\(k\).

For every formal \(k\)-group \(G\), denote by
\(\mathbf{Aut}_k(G)\) the \(k\)-group functor assigning to every
finite-dimensional commutative \(k\)-algebra \(C\) the group of
automorphisms of the formal \(C\)-group
\[
  G\widehat\otimes_k C.
\]
\footnote{\label{sga3:VIIB:nde:132}\textit{Editors' note.} The original stated:
``This \(k\)-functor is manifestly left exact
\textup{(\(H\) is topologically flat over \(k\)!)}.'' The editors
supplied the details that follow.}
Since \(k\) is a field, \(G\) is topologically flat over \(k\).
Equivalently, its affine algebra
\[
  A=\mathcal A(G)
\]
is topologically flat over \(k\), or its covariant bialgebra
\[
  H=\mathbf H(G)
\]
is flat over \(k\). The functor \(\mathbf{Aut}_k(G)\) is therefore
a formal \(k\)-group.

Indeed, \(\mathbf{Aut}_k(G)\) is the fiber product of the following
diagram \textup{(cf.\ \SGARefLink{sga3:i:remark:1-7-3}{Exposé~I, 1.7.3})}:
\begin{center}
  \phantomsection
  \label{sga3:VIIB:3.3.3:automorphism-fiber-product-diagram}
  \begin{tikzcd}[row sep=3.2em,column sep=3.8em]
{} &
\underline{\operatorname{End}}_k(G)
  \times\underline{\operatorname{End}}_k(G) \arrow[d] \\
\operatorname{Spf}(k) \arrow[r] &
\underline{\operatorname{End}}_k(G)
  \times\underline{\operatorname{End}}_k(G)
\end{tikzcd}

\end{center}
The vertical arrow sends
\[
  (\phi,\psi)\longmapsto(\phi\circ\psi,\psi\circ\phi),
\]
whereas the horizontal arrow is the unit section. It is therefore
enough to show that the \(k\)-functor \(\mathbf{End}_k(G)\) is
represented by an object of \(\operatorname{Vaf}/k\); by \SGARefLink{sga3:VIIB:1.1}{1.1} and
\SGARefLink{sga3:VIIB:0.4.2}{0.4.2}, this amounts to showing that it is left exact, or,
equivalently, that it commutes with fiber products of
\(k\)-algebras.

To give an element of \(\mathbf{End}_k(G)(C)\) is equivalent to
giving a \(k\)-algebra morphism
\[
  \phi:H\longrightarrow H\otimes_k C
\]
that also respects the coalgebra structure, that is, for which the
usual compatibility diagrams commute. Since \(H\) is flat over
\(k\), the functor
\[
  H\otimes_k -
\]
commutes with fiber products of \(k\)-algebras. Thus
\(\mathbf{End}_k(G)\) is left exact and can be treated as a formal
\(k\)-group.

If \(L\) is a \(k\)-Lie algebra and
\(G=\mathscr G(L)\), \SGARefLink{sga3:VIIB:3.3:cartier-theorem}{Theorem~3.3} shows that
\(\mathbf{Aut}_k(G)\) is isomorphic to the \(k\)-group functor
\(\mathbf{Aut}_k(L)\) which assigns to each finite-dimensional
\(k\)-algebra \(C\) the group of automorphisms of the
\(C\)-Lie algebra
\[
  C\otimes_k L.
\]

Now let \(G\) be an arbitrary formal \(k\)-group. We saw in \SGARefLink{sga3:VIIB:2.5.2:proposition}{2.5.2}
that \(G\) decomposes canonically as a semidirect product of an
étale formal group \(G_e\) and an infinitesimal formal group
\(G_0\). This semidirect product is determined by the data of
\[
  L=\operatorname{Lie}(G)=\operatorname{Lie}(G_0),
\]
the \(\Gamma\)-group
\[
  M=G_e(\overline k)=G(\overline k),
\]
and a morphism of formal \(k\)-groups
\[
  \Phi:
  G_e\longrightarrow
  \mathbf{Aut}_k(G_0)
  \simeq
  \mathbf{Aut}_k(L).
\]
Such a morphism sends \(G_e\) into the étale part
\(\mathbf{Aut}_k(L)_e\) of \(\mathbf{Aut}_k(L)\)
\textup{(cf.\ \SGARefLink{sga3:VIIB:2.5.1:corollary}{2.5.1})}. It is consequently determined by a morphism
of \(\Gamma\)-groups
\[
  \phi=\Phi(\overline k):
  M\longrightarrow
  \bigl(\mathbf{Aut}_kL\bigr)(\overline k)
  =
  \operatorname{Aut}_{\overline k}
    \bigl(L\otimes_k\overline k\bigr).
\]

Let \(\Gamma\) act on \(L\otimes_k\overline k\) by
\[
  \gamma(\ell\otimes t)=\ell\otimes\gamma(t).
\]
Then \(\phi\) makes \(M\) act on
\(L\otimes_k\overline k\) by automorphisms of
\(\overline k\)-Lie algebras in such a way that
\[
  \phi\bigl(\gamma(m)\bigr)
  =
  \gamma\circ\phi(m)\circ\gamma^{-1}.
\]
Equivalently,
\[
  \gamma(m)(\ell\otimes t)
  =
  \gamma\!\left(
    m\bigl(\ell\otimes\gamma^{-1}(t)\bigr)
  \right)
\]
for every \(m\in M\), \(\gamma\in\Gamma\), and
\(\ell\otimes t\in L\otimes_k\overline k\). We express this last
condition by saying that \(M\) acts on
\(L\otimes_k\overline k\) \emph{compatibly with \(\Gamma\)}.

The situation can be summarized in a ``highbrow'' statement. A
\emph{\(\Gamma\)-Lie algebra over \(k\)} is a triple
\[
  (L,M,\phi)
\]
consisting of a \(k\)-Lie algebra \(L\), a \(\Gamma\)-group \(M\),
and an action \(\phi\) of \(M\) on
\(L\otimes_k\overline k\) that is compatible with the actions of
\(\Gamma\) on \(M\) and on \(L\otimes_k\overline k\).

If \((L,M,\phi)\) and \((L',M',\phi')\) are two such
\(\Gamma\)-Lie algebras, a morphism from the first to the second is
a pair \((f,\theta)\), where
\[
  f:L\longrightarrow L'
\]
is a morphism of \(k\)-Lie algebras and
\[
  \theta:M\longrightarrow M'
\]
is a morphism of \(\Gamma\)-groups, such that
\[
  (f\otimes1)(m\cdot x)
  =
  \theta(m)\cdot(f\otimes1)(x)
\]
for every \(x\in L\otimes_k\overline k\) and \(m\in M\).
We have therefore proved:

\medskip
\noindent\textbf{Theorem.}
\phantomsection
\label{sga3:VIIB:3.3.3:gamma-lie-equivalence}
{\itshape Let \(k\) be a field of characteristic \(0\). The functor
that assigns to a formal \(k\)-group \(G\) the triple
\[
  \bigl(
    \operatorname{Lie}(G),
    G(\overline k),
    \Phi(\overline k)
  \bigr)
\]
is an equivalence from the category of formal \(k\)-groups to the
category of \(\Gamma\)-Lie algebras.}
